\begin{table}[htbp]
\centering
\caption{Cross-Dataset Comparison Across Three Real Submarine DAS Deployments}
\begin{tabular}{lrrr}
\toprule
Property & Dataset\_A\_Marlinks & Dataset\_B\_EMSO & Dataset\_C\_Oliktok \\
\midrule
Dataset Name & Marlinks Offshore Wind Farm DAS & EMSO Western Ionian Observatory DAS & Dryad Oliktok Submarine DAS \\
Geographic Location & Belgian North Sea (Offshore Export Cable) & Western Ionian Sea (Deep Seafloor, 2,100 m) & Beaufort Sea, Arctic Ocean (Oliktok Point) \\
Scientific Role & AIS-Corroborated Vessel Proximity \& CPA & Independent Seafloor Acoustic Baseline & Independent real-DAS environmental robustness \& false-escalation validation \\
Recording Duration & 590 s (10 min sequence) & 1,050 s (17.5 min recording) & 2,386,800 s (27.6 days continuous) \\
Spatial Channel Count & 250 channels (channels 1440-1690) & 2,963 channels along seafloor cable & 183 channels (indices 1088-3672) \\
Data Format & HDF5 (.h5) & NumPy (.npy) & NetCDF-4 / HDF5 (.nc) \\
Vessel / AIS Ground Truth & Continuous 1D vessel distance y (26 m - 2.8 km) & unavailable & unavailable \\
Independent Reference & Ship dimensions (304 m container ship) & INGV Catania deep-sea observatory & Seafloor oceanographic wave moorings \\
Temporal Stability (CV) & N/A (active transit, non-stationary) & 0.0008 & 0.1239 \\
Spatial Variability (CV) & Gini = 0.8739 (strong spatial peak at CPA) & CV = 1.7870 & CV = 0.2349 (Gini = 0.1265) \\
Primary Empirical Finding & Spearman rho = 0.9482 (p < 1e-30) vs 1/y & Extremely quiet acoustic baseline (kurtosis 13.2) & No T2/T3 vessel-escalation decisions across 27.6 days \\
\bottomrule
\end{tabular}
\end{table}
